\begin{claim}{D.67}{Proved}{The standard negative normalizers and their finite overgroups}
Retain the two standard negative involutions $b_0,b_1$
and the actual automorphisms $\alpha_j,\zeta_j$ of D.64.
Put
\[
 R(X)=Y,\quad R(Y)=X^{-1},\quad r=[R],\quad
 \delta=r^2=[\zeta_1],\quad
 V_j=C_\Omega(b_j),\quad D=\langle r,b_0\rangle .          \tag{D.67a}
\]
Then $D$ is the specified geometric dihedral group of order
eight, and
\[
 \begin{gathered}
 b_1=rb_0,\qquad
 V_0=\langle\delta,b_0\rangle,\qquad
 V_1=\langle\delta,b_1\rangle,\qquad
 V_0\cap V_1=\langle\delta\rangle,\\
 N_\Omega(V_0)=N_\Omega(V_1)=D,\qquad
 N_\Omega(D)=D,\qquad C_\Omega(D)=\langle\delta\rangle.
 \end{gathered}                                               \tag{D.67b}
\]
The complete list of finite subgroups of $\Omega$ containing
this specified $b_0$ is
\[
                     \langle b_0\rangle,\qquad V_0,\qquad D.
                                                               \tag{D.67c}
\]
Thus $D$ is its unique maximal finite overgroup. Conjugation
by any $\theta\in\Omega$ transports this exact list and all
the normalizer and centralizer equalities to the corresponding
specified conjugate groups. The statement does not assert
that every negative involution is conjugate to $b_0$ or $b_1$.
\end{claim}
\begin{proof}
\emph{Actual word maps for the dihedral group.}
Let $B(X)=X$, $B(Y)=Y^{-1}$. Its relation with the
cusp-normalized representative is
\[
                         \alpha_0=\operatorname{Inn}(Y)B.
                                                               \tag{D.67d}
\]
All of $R,B,\alpha_j,\zeta_j$ are actual automorphisms of
the free group and extend to its full completion. The
following identities are identities of automorphisms,
checked on the two free generators:
\[
 \begin{aligned}
 R^4&=1,& R^2&=\zeta_1,& BRB&=R^{-1},\\
 RB&=\alpha_1,& R\alpha_0&=\operatorname{Inn}(X^{-1})\alpha_1,\\
 R\alpha_0R^{-1}&=\operatorname{Inn}(X^{-1})\zeta_0,\\
 R\zeta_0R^{-1}&=\operatorname{Inn}(Y^{-1})\alpha_0,\\
 R\alpha_1R^{-1}&=\zeta_1\alpha_1,&
 \alpha_0\zeta_0&=\operatorname{Inn}(Y)\zeta_1.
 \end{aligned}                                               \tag{D.67e}
\]
Their peripheral maps, with $c_0=[X,Y]$, include
\[
 R(c_0)=X^{-1}c_0X,\qquad
 B(c_0)=Y^{-1}c_0^{-1}Y.                                  \tag{D.67f}
\]
Consequently the indicated outer classes belong to the
full peripheral group. The homology matrices of $R$ and
$B$ are respectively
\[
                 \begin{pmatrix}0&-1\\1&0\end{pmatrix},
                 \qquad\begin{pmatrix}1&0\\0&-1\end{pmatrix}.
                                                               \tag{D.67g}
\]
Their eight matrices $R_*^k B_*^\epsilon$, $0\leq k<4$,
$0\leq\epsilon<2$, are distinct. The actual relations
(D.67e) give at most eight elements, while these matrices
give at least eight distinct outer classes. Hence
$D\cong D_8$ with exactly these labels.
The linear maps $(u,v)\mapsto(-v,u)$ and
$(u,v)\mapsto(u,-v)$ on the torus, punctured at the origin,
realize these homology matrices. The discrete
once-punctured-torus identification
$\operatorname{Out}(F(X,Y))\cong\operatorname{GL}_2(\Z)$
proved in A.13 identifies their induced outer classes with
the displayed word maps. This supplies the asserted actual
geometric realization.

By D.64--D.65 the full centralizers are exactly
$V_0=\langle b_0,[\zeta_0]\rangle$ and
$V_1=\langle b_1,\delta\rangle$, both of order four.
Equations (D.67d)--(D.67e) give
$[\zeta_0]=\delta b_0$ and $b_1=rb_0$, proving the
first line of (D.67b) and identifying their intersection.
They are the two specified index-two Klein groups in $D$.

\emph{Their full normalizers.}
The unique nonidentity multiplier-one element of each
$V_j$ is $\delta$; its two other nonidentity elements have
multiplier $-1$. The multiplier homomorphism on $\Omega$
is conjugacy invariant. Thus conjugation by a member of
$N_\Omega(V_j)$ fixes $\delta$ and either fixes or exchanges
the other two elements. Its image in $\operatorname{Aut}(V_j)$
therefore has order at most two. Its kernel is
$C_\Omega(V_j)$, and the exact inclusions
\[
 V_j\leq C_\Omega(V_j)\leq C_\Omega(b_j)=V_j
\]
make this kernel precisely $V_j$. It follows that the full
normalizer has at most eight elements. The element $r$
realizes the nontrivial exchange in each case by (D.67e),
and the group $D$ already normalizes both $V_j$.
This proves both full normalizer equalities.

The mod-two homology map on $D$ has exact kernel $V_0$:
$B_*\equiv I$ and $R_*\equiv\begin{psmallmatrix}0&1\\1&0\end{psmallmatrix}$
over $\F_2$. Every element of $N_\Omega(D)$ preserves this
kernel, since its actual homology map conjugates the
mod-two matrices and preserves the identity matrix.
It therefore belongs to $N_\Omega(V_0)=D$.
Thus $N_\Omega(D)=D$. The full centralizer is now its
ordinary group center, which is $\langle\delta\rangle$.

\emph{Every finite overgroup of the specified two-class reflection.}
Let $K\leq\Omega$ be finite and contain $b_0$.
D.55 gives the complete nine-type list
$1,C_2,C_3,C_4,C_2^2,C_6,S_3,D_8,D_{12}$.
The first and third types cannot contain an involution.
The unique involution of a cyclic order-four group is
positive by D.55; that of a cyclic order-six group is
positive by D.56. Hence neither can be $b_0$.

We exclude $S_3$ and $D_{12}$ by the exact mod-two
representation. Every element $h$ of order three has
mod-two homology matrix of exact order three by D.54.
In $S_3$ every involution inverts its normal $C_3$.
In $D_{12}$ the central involution is the positive
involution of a cyclic-six subgroup, so a negative
involution is a reflection and also inverts the
normal $C_3$. Thus either type would give
\[
                    b_0 h b_0^{-1}=h^{-1},\qquad |h|=3.
                                                               \tag{D.67h}
\]
But $A_2(b_0)\bmod2=I$. Applying the actual homology
map and reducing modulo two would make the order-three
matrix equal to its inverse, a contradiction.

For $K\cong C_2^2$, every element commutes with $b_0$,
so $K\leq V_0$, and equality follows from their orders.
For $K\cong D_8$, its central involution is the square
of an order-four element, hence is positive. It
commutes with $b_0$, so it is the unique positive
nonidentity element $\delta$ of $V_0$.
The group $\langle\delta,b_0\rangle=V_0$ has index two
in $K$ and is normal. Consequently
$K\leq N_\Omega(V_0)=D$, and equality again follows
from the orders. The remaining cyclic-two group is
exactly $\langle b_0\rangle$. This proves (D.67c).

Finally $K\mapsto\theta K\theta^{-1}$ and its inverse
are exact bijections between finite overgroups of $b_0$
and of $\theta b_0\theta^{-1}$. The same conjugation
identifies each full normalizer and centralizer.
This gives the asserted transport without any claim
that these specified orbits exhaust all negative actions.

The executable certificate
\path{code/part_c/standard_negative_overgroups_d67.py}
checks the actual word identities, all eight signed
permutation automorphisms and their multiplication,
both labelled Klein subgroups, and the mod-two
order-three inversion obstruction. The full upper
bounds use the proved centralizers and nine-type list,
not a finite enumeration of the ambient profinite group.
\end{proof}
